\subsection{Fonctions Récursives}

Une fonction récursive est une fonction qui \textit{s'appelle elle-même} à l'intérieur de son propre bloc d'instructions.

\begin{verbatim}
def f_r([parametres]):
  # 1. Condition d'arrêt
  if [condition_de_fin_atteinte]:
     return [valeur_finale]
    
  # 2. Appel Récursif 
  else:
    # Modification possible
    [new_p] = [modification]
        
# La fonction s'invoque elle-même
    return [operation] + f_r([new_p])
\end{verbatim}

Les composants indispensables sont :

\begin{itemize}
\item Condition d'arrêt : Évalue si le problème est suffisamment réduit pour être résolu directement sans nouvel appel
\item Modification d'état : Les arguments passés à l'appel suivant doivent systématiquement se rapprocher de la condition d'arrêt
\item Le mot-clé \texttt{return} : Permet de faire remonter le résultat d'un appel au précédent lors du "dépilage".
\end{itemize}